\documentclass[10pt,a4paper]{amsart}
\usepackage[margin=19mm]{geometry}
\usepackage{amsmath,amssymb,mathtools}
\usepackage[scr=rsfs,cal=euler]{mathalfa}
\usepackage{xfrac}
\usepackage[unicode,hidelinks,bookmarks=false]{hyperref}
\newcommand{\ie}{i.e.\ }
\newcommand{\cf}{cf.\ }
\newcommand{\resp}{respectively\ }
\newcommand{\Subsection}{\subsection}
\providecommand{\nobreakdash}{\nobreak\hbox{-}}
\setlength{\emergencystretch}{2em}
\multlinegap0pt
\usepackage{amssymb}
\usepackage{algorithm}\usepackage{algpseudocode}
\usepackage{enumerate}
\usepackage{float}
\usepackage{mathtools}
\usepackage{tcolorbox}
\usepackage{bm}
\newtheorem{assumption}{Assumption}
\newtheorem{lemma}{Lemma}
\newtheorem{theorem}{Theorem}
\renewcommand{\Re}{{\mathbb{R}}}
\title{\LARGE \bf A Unifying Primal-Dual Proximal Framework for Distributed Nonconvex Optimization}
\author{Zichong Ou, Jie Lu}
\date{Source: arXiv:2603.09524v1 (CC BY 4.0)}
\begin{document}
\maketitle

\noindent\emph{Source and licence.} \LARGE \bf A Unifying Primal-Dual Proximal Framework for Distributed Nonconvex Optimization. Version arXiv:2603.09524v1 (CC BY 4.0), distributed by arXiv under the Creative Commons Attribution 4.0 International licence, http://creativecommons.org/licenses/by/4.0/, as recorded in the arXiv licence metadata for this exact version and shown on the abstract page https://arxiv.org/abs/2603.09524v1. Full licence notice: \texttt{licenses/arxiv-2603.09524-CC-BY-4.0.txt}.\par\smallskip

\noindent\textit{Editorial scope.} Counted unit: the article's theorem Bk (source0.tex lines 823--845). The article's complete original proof of it is delivered with it (source0.tex lines 846--848, one \texttt{proof} environment of 3 lines). No mathematics is written, repaired or re-derived: the statement and the proof are the article's own lines, in the article's own notation, and no retained line is substituted or re-flowed.  Retained uncounted as the source-local context the proof needs: lines 120--122; lines 204--209; lines 211--213; lines 289--297; lines 745--747; lines 749--751; lines 761--767; lines 795--798; lines 800--802; lines 810--816; lines 1601--1603; lines 1624--1671.  Nothing was omitted from any retained span, so the package carries no omission record.  No retained line contains a citation, so the wrapper carries no bibliography and no bibliography entry.  No retained line contains a figure environment, an image-inclusion command or drawing code, so no image is delivered and the package declares no asset dependency.  Editorial formatting only, in addition: an AMS \texttt{amsart} preamble that restates the article's own declarations and macros used by the retained lines, namely the environments \texttt{assumption}, \texttt{lemma}, \texttt{theorem} on separate counters, together with the standard packages those lines need.  The retained environments print in this document as Assumption 1, Lemma 1, Theorem 1 in order of appearance, whereas the article prints the same items with the numbering of its own full document.  The complete native source of the article as distributed in the arXiv e-print for this exact version is bundled unchanged as \texttt{sources/196070/source0.tex}.
\begin{equation}
\min _{x \in \mathbb{R}^{d}} f(x)= \sum_{i=1}^{N} f_{i}(x), \label{p1}
\end{equation}

\begin{assumption}\label{assumption smooth}
The global objective $\tilde{f}:\Re^{Nd}\rightarrow\Re$ is $\bar{M}$-smooth for some $\bar{M}> 0$, i.e.,
\begin{equation}
\|\nabla \tilde{f}(\mathbf{x}) - \nabla \tilde{f}(\mathbf{y})\|\leq \bar{M}\|\mathbf{x}-\mathbf{y}\|,\quad\forall \mathbf{x},\mathbf{y}\in \Re^{Nd}.\label{smooth}
\end{equation}
\end{assumption}

\begin{assumption}\label{assumption: optimality}
The optimal set $\mathbb{X}^*$ of problem~\eqref{p1} is nonempty and the optimal value $f^* >-\infty$.
\end{assumption}

\begin{assumption}\label{assumption: distributed}
    The matrix $\mathbf{P}=(p_{ij})_{N\times N}$ satisfies the following conditions:
    \begin{itemize}
        \item[($\romannumeral1$)] $\mathbf{P}$ is symmetric and positive semi-definite.
        \item[($\romannumeral2$)] $\operatorname{Null}(\mathbf{P})=\operatorname{span}(\mathbf{1}_N)$.
        \item[($\romannumeral3$)] $p_{ij}=p_{ji}>0, \forall \{i,j\}\in \mathcal{E}$.
        \item[($\romannumeral4$)] $p_{ij}=p_{ji}=0, \forall i\in \mathcal{V},\,\forall j\notin \mathcal{N}_i\cup \{i\}$.
    \end{itemize}
\end{assumption}

	\begin{equation}
		\label{first-order opt modified} \nabla \tilde{f}(\mathbf{x}^k) + \mathbf{L}^{\frac{1}{2}}\mathbf{v}^k + \rho \mathbf{L} \mathbf{x}^{k+1}+ \mathbf{B}^k(\mathbf{x}^{k+1} - \mathbf{x}^k)=0,
	\end{equation}

	\begin{equation}
		\mathbf{v}^{k+1}=\mathbf{v}^k+\rho \mathbf{L}^{\frac{1}{2}}\mathbf{x}^{k+1}. \label{vk+1 opt}
	\end{equation}

\begin{lemma} \label{lemma vk+1 - vk}
	Suppose Assumptions~\ref{assumption smooth}--\ref{assumption: distributed} hold. Let $\{\mathbf{x}^k\}$ be the sequence generated by UPP-SC with $\mathbf{O}\prec \lambda_N^{\mathbf{B}^k} \mathbf{I}\preceq \mathbf{B}^{k}\preceq \lambda_1^{\mathbf{B}^k}\mathbf{I} $. Then, for all $k\geq 0$,
	\begin{align}
		&\rho\|\mathbf{x}^{k+1}\|^2_{\mathbf{L}}=\frac{1}{\rho}\|\mathbf{v}^{k+1}-\mathbf{v}^k\|^2\leq 3\tilde{\kappa} \big((c_{\mathbf{B}}+\bar{M}^2)\|\mathbf{x}^k-\mathbf{x}^{k-1}\|^2_{(\mathbf{B}^{k})^{-1}}+\|\mathbf{w}^{k+1}\|^2_{\mathbf{B}^k}\big), \label{vk+1 - vk}
	\end{align}
	where $c_{\mathbf{B}}=4(\lambda_1^\mathbf{B})^2$ and $\{\mathbf{v}^k\}_{k=0}^\infty$ is given by \eqref{vk+1 opt}.
\end{lemma}

	\begin{align}
		\label{define tilde c}& 1 \geq \tilde{c}\geq 9\tilde{\kappa}=\max_{k\geq 0}\frac{9\lambda_1^{\mathbf{B}^k}}{\rho \lambda_{N-1}^{\mathbf{L}}}>0, \\
		\label{Bk+rhoL-c/2Bk-1}& \frac{1}{4}(\mathbf{B}^k+\rho \mathbf{L})-\frac{\tilde{c}}{2}(\mathbf{B}^{k+1}-\mathbf{B}^k+\mathbf{B}^{k-1})\!-\!(\frac{1}{2}+2\tilde{c})\bar{M}\mathbf{I}_{Nd}-\frac{\tilde{c}(2+\tilde{c})}{6}(c_{\mathbf{B}}+\bar{M}^2)(\mathbf{B}^{k+1})^{-1} \succeq \mathbf{O}_{Nd}.
	\end{align}

	\begin{align}
		\tilde{P}^{k+1}-\tilde{P}^{k}\leq& -\|\mathbf{x}^{k+1}-\mathbf{x}^k\|^2_{\frac{1}{4}(\mathbf{B}^k+\rho \mathbf{L})}-\frac{\tilde{c}}{2} \|\mathbf{w}^{k+1}\|^2_{\frac{(1-\tilde{c})}{3}\mathbf{B}^k+\mathbf{B}^{k-1}}\leq 0.  \label{tildePk+1 - tildePk lemma}
	\end{align}

\begin{lemma}\label{lemma tildePk+1 geq f*}
	Suppose all the conditions in Lemma~\ref{lemma vk+1 - vk} hold. Let $\{\mathbf{x}^k\}$ be the sequence generated by UPP-SC with the parameters selected by \eqref{define tilde c} and \eqref{Bk+rhoL-c/2Bk-1}. Also let $\mathbf{x}^{-1}=\mathbf{v}^{-1}=0$, and $\mathbf{G}^{-1}=\mathbf{G}^0$. Then,
	\begin{equation*}
		\tilde{P}^{k+1}\geq f^*, \forall k >0, \quad \tilde{P}^{0}\leq \tilde{f}(\mathbf{x}^0)+\frac{6\tilde{c}+7}{8\bar{M}}\|\nabla \tilde{f}(0)\|^2,
	\end{equation*}
	where $f^*$ is defined in Assumption~\ref{assumption: optimality}.
\end{lemma}

\begin{equation}
	\hat{P}^{k+1}\geq 0, \quad \tilde{P}^{k+1}\geq f^*, \quad \forall k \geq 0. \label{hat Pk+1 geq 0 tildePk+1 geq f*}
\end{equation}

\begin{align}
	\tilde{P}^0\leq& \tilde{f}(\mathbf{x}^0)+ (\frac{3}{2}\tilde{c}+\frac{7}{4})\|\nabla \tilde{f}(0)\|^2_{\mathbf{G}^0}\notag\\
	\overset{\eqref{Bk+rhoL-c/2Bk-1}}{\leq} &\tilde{f}(\mathbf{x}^0) +\frac{(\frac{3}{2}\tilde{c}+\frac{7}{4})(\frac{1}{4}+\frac{\tilde{c}}{2})}{(\frac{1}{2}+2\tilde{c})\bar{M}}\|\nabla \tilde{f}(0)\|^2 \notag\\
	< & \tilde{f}(\mathbf{x}^0) +\frac{6\tilde{c}+7}{8\bar{M}}\|\nabla \tilde{f}(0)\|^2. \label{tildeP0 leq fx0+nabla f0}
\end{align}
Hence, from \eqref{hat Pk+1 geq 0 tildePk+1 geq f*} and \eqref{tildeP0 leq fx0+nabla f0}, we obtain Lemma~\ref{lemma tildePk+1 geq f*}.
\subsubsection{Proof of Theorem~\ref{theorem UPP-SC sublinear time-varying Bk}}\label{proof theorem UPP-SC sublinear time-varying Bk}

\textbf{(i)} First we illustrate the feasibility of the parameters in \eqref{define tilde c} and \eqref{Bk+rhoL-c/2Bk-1}. 

% To show \eqref{Bk+rhoL-c/2Bk-1}, we consider its feasibility in the range space of $\mathcal{S}$ and its complement space $\mathcal{S}^{\bot}$, respectively.

% For the former, we multiply $\mathbf{J}\in \mathcal{S}$ to the both sides of \eqref{Bk+rhoL-c/2Bk-1} and have
% \begin{equation}
% 	\label{condition lambdassumption: optimality} \frac{1}{4}\underline{\lambda}_1-\tilde{c}\bar{\lambda}_1-(\frac{1}{2}+2\tilde{c})\bar{M}-\frac{\tilde{c}(2+\tilde{c})}{6\underline{\lambda}_1}\bar{M}^2>0.
% \end{equation}
% From \eqref{tildec} and \eqref{lambdaNB range}, we know that the inequality~\eqref{condition lambdassumption: optimality} holds. 

% Then, we show the feasibility of \eqref{Bk+rhoL-c/2Bk-1} in the complement space $\mathcal{S}^{\bot}$. From \eqref{define tilde c}, it is sufficient to show that
% \begin{equation}
% 	\frac{1}{4}(\rho \underline{\lambda}_\mathbf{L}+\underline{\lambda}_2)-\tilde{c}d_2\underline{\lambda}_2-(\frac{1}{2}+2\tilde{c})\bar{M}-\frac{\tilde{c}(2+\tilde{c})\bar{M}^2}{6\underline{\lambda}_2}>0. \label{condition rho lambdaL}
% \end{equation}
To see \eqref{Bk+rhoL-c/2Bk-1}, it is sufficient to show that $\frac{1}{4}\lambda_N^{\mathbf{B}}-\frac{\tilde{c}}{2}(\lambda_1^{\mathbf{B}}-\lambda_N^{\mathbf{B}}+\lambda_1^{\mathbf{B}})-(\frac{1}{2}+2\tilde{c})\bar{M}
-\frac{\tilde{c}(2+\tilde{c})}{6}(c_\mathbf{B}+\bar{M}^2)/\lambda_N^{\mathbf{B}}> 0,$
% \begin{align*}
% 	&\frac{1}{4}\lambda_N^{\mathbf{B}}-\frac{\tilde{c}}{2}(\lambda_1^{\mathbf{B}}-\lambda_N^{\mathbf{B}}+\lambda_1^{\mathbf{B}})-(\frac{1}{2}+2\tilde{c})\bar{M}\notag\\
% 	& -\frac{\tilde{c}(2+\tilde{c})}{6}(c_\mathbf{B}+\bar{M}^2)/\lambda_N^{\mathbf{B}}> 0,
% \end{align*}
which can be transformed into the following inequality:
\begin{equation}\label{d1lambdaNb-d2-d3/lambdaNB}
	d_1\lambda_N^{\mathbf{B}}-d_2-d_3/\lambda_N^{\mathbf{B}} >0,
\end{equation}
where $d_1= \frac{1}{4}- \xi_5 \tilde{c}-\xi_6 \tilde{c}^2$, $d_2 = (\frac{1}{2}+2\tilde{c})\bar{M}$, and $d_3={(2+\tilde{c})\tilde{c}}\bar{M}^2/6$.
Since $0<\tilde{c}<({-\xi_5+\sqrt{\xi_5^2+\xi_6}})/({2\xi_6})$ implies $d_1>0$, together with \eqref{d1lambdaNb-d2-d3/lambdaNB}, we obtain $\lambda_N^{\mathbf{B}}\geq ({d_2+\sqrt{d_2^2+4d_1d_3}})/({2d_1})$. Hence, \eqref{Bk+rhoL-c/2Bk-1} is satisfied. Besides, from $\rho \geq ({9\kappa_{\mathbf{B}}\lambda_N^{\mathbf{B}}})/({\tilde{c}\lambda_{N-1}^{\mathbf{L}}})$, we have \eqref{define tilde c}.

\textbf{(ii)}
Next, we analyze the convergence result of UPP-SC. 

By multiplying both sides of \eqref{first-order opt modified} by matrix $\mathbf{J}$, since $\mathbf{J}\mathbf{L}=0$ we obtain $\mathbf{J}\nabla \tilde{f}(\mathbf{x}^{k}) + \mathbf{J}\mathbf{B}^k(\mathbf{x}^{k+1}-\mathbf{x}^{k})=0$, for $k \geq -1.$
% \begin{equation*}
% 	\mathbf{J}\nabla \tilde{f}(\mathbf{x}^{k}) + \mathbf{J}\mathbf{B}^k(\mathbf{x}^{k+1}-\mathbf{x}^{k})=0,\quad \forall k \geq -1.
% \end{equation*}
From $\mathbf{J}\nabla \tilde{f}(\mathbf{x})=\mathbf{1}_N \otimes\big(\frac{1}{N} \sum_{i=1}^{N}\nabla\tilde{f}_i(x_i)\big)$, $\mathbf{J}\mathbf{J}=\mathbf{J}$, $\|\mathbf{J}\|=1$,
\begin{align}
	&\frac{1}{N}\|\sum_{i=1}^N \nabla \tilde{f}_i(x_i^k)\|^2 \leq (\mathbf{x}^{k+1}-\mathbf{x}^{k})^{\mathsf{T}}\mathbf{B}^k\mathbf{J}\mathbf{J}\mathbf{B}^k(\mathbf{x}^{k+1}-\mathbf{x}^{k}) \notag\\
	\leq & \|\mathbf{x}^{k+1}-\mathbf{x}^{k}\|^2_{\mathbf{B}^k}\|\mathbf{B}^k\| 
	\overset{\eqref{tildePk+1 - tildePk lemma}}{\leq} 4\|\mathbf{B}^k\|(\tilde{P}^k-\tilde{P}^{k+1}). \label{bounded gradient norm Bk}
\end{align}

\begin{theorem}\label{theorem UPP-SC sublinear time-varying Bk}
	Suppose all the conditions in Lemma~\ref{lemma tildePk+1 geq f*} hold. Let $\{\mathbf{x}^k\}$ be the sequence generated by UPP-SC with the parameters selected as follows: $\xi_5=\frac{4\kappa_{\mathbf{B}}^2}{3}+\frac{2\kappa_{\mathbf{B}}-1}{2}$, $\xi_6= \frac{2\kappa_{\mathbf{B}}^2}{3}$, $0<\tilde{c}<\frac{-\xi_5+\sqrt{\xi_5^2+\xi_6}}{2\xi_6}$, $d_1= \frac{1}{4}- \xi_5 \tilde{c}-\xi_6 \tilde{c}^2$, $d_2 = (\frac{1}{2}+2\tilde{c})\bar{M}$, $d_3=\frac{(2+\tilde{c})\tilde{c}}{6}\bar{M}^2$, $\lambda_1^{\mathbf{B}}\geq\lambda_N^{\mathbf{B}}\geq \frac{d_2+\sqrt{d_2^2+4d_1d_3}}{2d_1}$ and $\rho \geq \frac{9\kappa_{\mathbf{B}}\lambda_N^{\mathbf{B}}}{\tilde{c}\lambda_{N-1}^{\mathbf{L}}}$.
	% \begin{align}
	% 	&0<\tilde{c}<\frac{-\xi_5+\sqrt{\xi_5^2+\xi_6}}{2\xi_6}, \text{ with} \notag \\
	% 	&\quad\xi_5=\frac{4\kappa_{\mathbf{B}}^2}{3}+\frac{2\kappa_{\mathbf{B}}-1}{2}, \quad \xi_6= \frac{2\kappa_{\mathbf{B}}^2}{3},\label{tildec}\\
	% 	&\lambda_N^{\mathbf{B}}\geq \frac{d_2+\sqrt{d_2^2+4d_1d_3}}{2d_1} , \text{ with } d_1= \frac{1}{4}- \xi_5 \tilde{c}-\xi_6 \tilde{c}^2,\notag\\
	% 	& \quad d_2 = (\frac{1}{2}+2\tilde{c})\bar{M},\quad d_3=\frac{(2+\tilde{c})\tilde{c}}{6}\bar{M}^2, \label{lambdaNB range} \\
	% 	\label{rho range} &\rho \geq \frac{9\kappa_{\mathbf{B}}\lambda_N^{\mathbf{B}}}{\tilde{c}\lambda_{N-1}^{\mathbf{L}}}.
	% \end{align}
	Let $T_k$ denote the iteration index in which UPP-SC satisfies
	\begin{equation}
		e(T_k):= \min_{k\in[T_k]}\frac{1}{N}\|\sum_{i=1}^N \nabla {f}_i(x_i^k)\|^2+\rho\|\mathbf{x}^{k}\|^2_\mathbf{L}\leq \epsilon, \label{eT time-varying Bk}
	\end{equation}
	where the error $\epsilon>0$. Then, the error above has the following bound:
	\begin{equation}
		e(T_k) \leq \tilde{C}_1\times {\tilde{C}_2}/{T_k}, \label{epsilon leq tildeC1C2/T}
	\end{equation}
	where $\tilde{C}_1:= \tilde{f}(\mathbf{x}^0)-f^*+\frac{6\tilde{c}+7}{8\bar{M}}\|\nabla \tilde{f}(0)\|^2$ and $\tilde{C}_2:= 4{\lambda}_1^{\mathbf{B}}+\frac{4}{3(1-\tilde{c})}(1+\frac{18d_1}{\tilde{c}(\tilde{c}+2)})+8(\frac{6d_1}{\tilde{c}(\tilde{c}+2)}+1).$
	% \begin{align*}
	% 	\tilde{C}_1:= & \tilde{f}(\mathbf{x}^0)-f^*+\frac{6\tilde{c}+7}{8\bar{M}}\|\nabla \tilde{f}(0)\|^2,\\
	% 	\tilde{C}_2:= & 4{\lambda}_1^{\mathbf{B}}+\frac{4}{3(1-\tilde{c})}(1+\frac{18d_1}{\tilde{c}(\tilde{c}+2)})+8(\frac{6d_1}{\tilde{c}(\tilde{c}+2)}+1).
	% \end{align*}
\end{theorem}

\begin{proof}
    See Appendix~\ref{proof theorem UPP-SC sublinear time-varying Bk}.
\end{proof}

\end{document}
